\documentclass{article}
\usepackage[T1]{fontenc}
\usepackage{lmodern}
\usepackage{graphicx}
\usepackage{amsmath,amssymb}
\usepackage[a4paper,margin=2cm]{geometry}
\usepackage[absolute,overlay]{textpos}
\setlength{\TPHorizModule}{1mm}
\setlength{\TPVertModule}{1mm}
\usepackage[indonesian]{babel}
\usepackage{hyperref}
\setlength{\emergencystretch}{2em}
% unit: Halaman 24

\begin{document}

% === Halaman 24 (llm) ===

\section*{Kriptanalisis Caesar Cipher}

\begin{itemize}
    \item \textit{Caesar cipher} mudah dipecahkan dengan \textit{exhaustive key search} (\textit{brute force}) karena jumlah kuncinya sangat sedikit (hanya ada 26 kunci).
    \item Coba lakukan dekripsi dengan berbagai nilai $k$ dari 0 sampai 25, lalu periksa apakah hasil dekripsi merupakan kata atau kalimat yang bermakna. Jika ya, maka diduga $k$ adalah kuncinya.
    \item Untuk memastikan $k$ adalah kunci yang benar, maka cobakan $k$ untuk potongan kriptogram lainnya.
\end{itemize}

\end{document}
